\section*{第 \ref{ch:g12:mountain-building} 章——造山运动}

\begin{solution}{exo:g12:mountain-building:1}
(a) \omterm{def:g12:mountain-building:thrust}{逆冲}；(b) \omterm{def:g12:mountain-building:ophiolite}{蛇绿岩}；(c) \omterm{def:g12:mountain-building:fold}{向斜}；(d) \omterm{def:g12:mountain-building:suture}{缝合带}；(e) \omterm{def:g12:mountain-building:fold}{背斜}；
(f) \omterm{def:g12:mountain-building:nappe}{推覆体}。
\end{solution}

\begin{solution}{exo:g12:mountain-building:2}
(1) 两翼岩层从\omterm{def:g12:mountain-building:fold}{褶皱轴}向外倾斜：为\omterm{def:g12:mountain-building:fold}{背斜}。(2) 最老的
岩层。(3) 不对：侵蚀沿软硬岩层进行，而非沿\omterm{def:g12:mountain-building:fold}{褶皱}的形态进行；
\omterm{def:g12:mountain-building:fold}{背斜}的核部可能被侵蚀成山谷。
\end{solution}

\begin{solution}{exo:g12:mountain-building:3}
(1) $e = 1 - \cos\ang{30} = 1 - 0.866 = 0.134$，约 \qty{13}{\percent}。
(2) $L_0 = 6.0/\cos\ang{30} \approx \qty{6.9}{km}$。
\end{solution}

\begin{solution}{exo:g12:mountain-building:4}
(1) \omterm{def:g11:ocean-ridges:normal-fault}{正断层}；地壳发生了拉张。(2) 是的：这是一条\omterm{def:g12:mountain-building:thrust}{逆断层}，
倾角小于 \ang{45}。
\end{solution}

\begin{solution}{exo:g12:mountain-building:5}
(1) 海相沉积物、枕状玄武岩、席状\omterm{def:g11:igneous-rocks:intrusion}{岩墙}、\omterm{def:g11:ocean-ridges:gabbro}{辉长岩}、\omterm{def:g11:ocean-ridges:peridotite}{橄榄岩}。
(2) 熔岩只有在水中冷却时才会形成枕状体。(3) \omterm{def:g12:mountain-building:ophiolite}{蛇绿岩}是曾位于两大陆之间的一段\omterm{def:g11:ocean-ridges:oceanic-lithosphere}{大洋岩石圈}；
两\omterm{def:g12:mountain-building:collision}{大陆碰撞}时，其洋底\omterm{def:g11:minerals:thin-section}{薄片}被推挤到大陆边缘，
而未\omterm{def:g11:subduction-zones:subduction}{俯冲}下沉，因此它们
出现在大洋闭合之处，沿\omterm{def:g12:mountain-building:suture}{缝合带}分布。
\end{solution}

\begin{solution}{exo:g12:mountain-building:6}
(1) $L_0 = 2.1 + 1.8 + 2.4 + 1.6 = \qty{7.9}{km}$；
$e = (7.9 - 6.0)/7.9 \approx 0.24$，约 \qty{24}{\percent}。(2) $\cos\theta =
6.0/7.9 \approx 0.759$，$\theta \approx \ang{41}$。
\end{solution}

\begin{solution}{exo:g12:mountain-building:7}
(1) $45\cos\ang{10} \approx \qty{44.3}{km}$。(2) $45\sin\ang{10} \approx
\qty{7.8}{km}$。(3) 这类\omterm{def:g12:mountain-building:fold}{褶皱}的\omterm{def:g12:mountain-building:shortening}{缩短量}为 $1 - \cos\ang{12} \approx 0.0219$：
\omterm{def:g12:mountain-building:fold}{褶皱}之前，造山带的宽度必定为
$44.3/0.0219 \approx \qty{2000}{km}$。
\end{solution}

\begin{solution}{exo:g12:mountain-building:8}
(1) 叠置律：若忽略断层，上方的岩层
会被误认为较年轻。(2) $500 - 75 = 425$ 百万年。(3) $\qty{40000}{m} /
\qty{800}{m} = 50$ 倍。
\end{solution}

\begin{solution}{exo:g12:mountain-building:9}
(1) $(30.23 - 19.38) \times 111.2 = 10.85 \times 111.2 \approx
\qty{1207}{km}$。(2) $1207/14 \approx \qty{86}{km}$ 每百万年，
即 \qty{86}{mm}，或每年约 \qty{8.6}{cm}。(3) 每世纪 \qty{9}{m} 即
每年 \qty{9}{cm}：两者在约 \qty{5}{\percent} 的范围内一致。
\end{solution}

\begin{solution}{exo:g12:mountain-building:10}
(1) $H = 35/(1 - 0.40) \approx \qty{58}{km}$。(2) $\Delta H \approx
\qty{23.3}{km}$；它高出的部分为 $23.3 \times 550/3380 \approx
\qty{3.8}{km}$，其额外的\omterm{def:g12:continental-crust:root}{山根}深为 $23.3 - 3.8 \approx \qty{19.5}{km}$。
\end{solution}

\begin{solution}{exo:g12:mountain-building:11}
(1) 图中将 $\tan\theta$ 放大了 3 倍：$\tan\theta = \tan\ang{40}/3
\approx 0.280$，$\theta \approx \ang{15.6}$。(2) 正确：$1 - \cos\ang{15.6}
\approx 0.037$，即 \qty{3.7}{\percent}；按图读取：$1 - \cos\ang{40}
\approx 0.234$，即 \qty{23}{\percent}。(3) 竖直方向拉伸的比例尺会使倾角变陡，
并使岩层变长，因此使\omterm{def:g12:mountain-building:shortening}{缩短量}看起来
大了好几倍。
\end{solution}

\begin{solution}{exo:g12:mountain-building:12}
(1) $L_0 = 10/(1 - 0.317) \approx \qty{14.6}{km}$。(2) $1 - \cos\ang{57.5}
\approx 1 - 0.537 = 0.463$，即 \qty{46}{\percent}。(3) 在圆顶状\omterm{def:g12:mountain-building:fold}{褶皱}中，
岩层只有一小段以 \ang{57.5} 倾斜；靠近转折端处岩层几乎水平，
而水平部分不会被缩短。在直翼\omterm{def:g12:mountain-building:fold}{褶皱}中，
整层岩层都以 \ang{57.5} 倾斜。
\end{solution}

\begin{solution}{exo:g12:mountain-building:13}
(1) $\sqrt{16.6^2 + 41.8^2} \approx \qty{45.0}{mm/\text{年}}$。(2) $41.8 - 23.5 =
\qty{18.3}{mm}$ 每年，即拉萨以北的南北向缩短约
\qty{18}{km}（每百万年）。(3) 拉萨本身已向北移动，而造山带的地震
广泛延伸至喜马拉雅以北、遍及亚洲：
该处地壳至今仍在断裂。
\end{solution}

\begin{solution}{exo:g12:mountain-building:14}
(1) 其倾角接近于零，因此\omterm{def:g12:mountain-building:shortening}{缩短量}即为位移，
\qty{45}{km}。(2) $45/(1 - \cos\ang{12}) = 45/0.0219 \approx \qty{2060}{km}$。
(3) 大幅度的汇聚主要由\omterm{def:g12:mountain-building:thrust}{逆冲}、\omterm{def:g12:mountain-building:nappe}{推覆体}和紧闭\omterm{def:g12:mountain-building:fold}{褶皱}吸收；
平缓的\omterm{def:g12:mountain-building:fold}{褶皱}吸收的量很少。
\end{solution}

\begin{solution}{exo:g12:mountain-building:15}
(1) $\qty{8848.86}{m} = \qty{8848860}{mm}$，历时 $\num{5e7}$ 年：平均每年约
\qty{0.18}{mm}。(2) $40/0.18 \approx 220$：隆升速度比印度的运动
慢二百多倍。(3) 石灰岩沉积于海平面以下，
其\omterm{def:g11:fossil-fuel-geology:petroleum-system}{上覆岩层}后被侵蚀殆尽：它被抬升的
幅度已超过其现今的海拔高度。
\end{solution}

\begin{solution}{pb:g12:mountain-building:1}
\textbf{1.} $2\pi \times 6371/360 \approx \qty{111.2}{km}$。

\textbf{2.} 在 5000 万年前（南纬 \ang{0.68}）与 4500 万年前（北纬 \ang{3.30}）之间：
$50 - 5 \times 0.68/3.98 \approx 49$ 百万年前。

\textbf{3.} $(19.38 - 0.68) \times 111.2 = 18.70 \times 111.2 \approx
\qty{2079}{km}$。

\textbf{4.} $2079/16 \approx \qty{130}{km}$ 每百万年，即 \qty{130}{mm}
每年，约 \qty{13}{cm}。

\textbf{5.} 66 至 60：$7.28 \times 111.2/6 \approx 135$；60 至 55：$4.86
\times 111.2/5 \approx 108$；55 至 50：$6.56 \times 111.2/5 \approx 146$；50
至 45：$3.98 \times 111.2/5 \approx 89$；45 至 40：$2.29 \times 111.2/5
\approx 51$（\unit{mm} 每年）。速度最快的是 5500 万至 5000 万年前。

\textbf{6.} $(20.00 - 5.59) \times 111.2/40 = 14.41 \times 111.2/40 \approx
\qty{40}{mm}$ 每年。

\textbf{7.} $40/130 \approx 0.31$：不到三分之一。

\textbf{8.} 从 55--50 阶段（146）到 50--45 阶段（89）的
转变大约发生在 5000 万年前。

\textbf{9.} \omterm{def:g12:continental-crust:continental-crust}{大陆地壳}太轻，无法下沉：当它抵达
海沟时，\omterm{def:g11:subduction-zones:subduction}{俯冲}已无法再自由地牵引印度，另一块
大陆对它起到了阻挡作用。

\textbf{10.} $(20.00 + 0.68) \times 111.2 = 20.68 \times 111.2 \approx
\qty{2300}{km}$.

\textbf{11.} $(29.65 - 21.15) \times 111.2 = 8.50 \times 111.2 \approx
\qty{945}{km}$.

\textbf{12.} $2300 - 945 \approx \qty{1354}{km}$.

\textbf{13.} 今日 $(29.65 - 20.00) \times 111.2 \approx \qty{1073}{km}$；
5000 万年前 $(21.15 + 0.68) \times 111.2 = 21.83 \times 111.2
\approx \qty{2427}{km}$（且 $2427 - 1073 = 1354$）。

\textbf{14.} $e = 1354/2427 \approx 0.56$，\qty{56}{\percent}。

\textbf{15.} \omterm{def:g12:mountain-building:fold}{褶皱}，以及\omterm{def:g12:mountain-building:thrust}{逆断层}：即携带\omterm{def:g12:mountain-building:nappe}{推覆体}的\omterm{def:g12:mountain-building:thrust}{逆冲}。

\textbf{16.} $H = 35 \times 2427/1073 \approx \qty{79}{km}$.

\textbf{17.} $\Delta H = 79 - 35 = \qty{44}{km}$；它将
高出 $44 \times 550/3380 \approx \qty{7.2}{km}$。

\textbf{18.} $\Delta H = 4.6 \times 3380/550 \approx \qty{28}{km}$，地壳
厚度为 $35 + 28 = \qty{63}{km}$，接近
实测的 \qty{70}{km}。

\textbf{19.} 地壳的截面并不保持不变：亚洲的部分地区被
向东挤出，同时侵蚀作用从
山顶移走岩石。

\textbf{20.} $2079 + 2300 \approx \qty{4380}{km}$：自白垩纪末以来，
印度已大约\textbf{向北移动 \qty{4400}{km}}。
\end{solution}
